% Use this file to enter your text.
\section{Introduction}

People routinely use music to maintain or change how they feel, and the
emotional fit of a track is a strong driver of listening
satisfaction~\cite{juslin2008}. Streaming services recommend at enormous scale,
but most recommenders optimise for long-run engagement inferred from listening
history rather than for the listener's \emph{current} affective state. As a
result, the same playlist is served whether a listener is tense and restless or
calm and content.

A growing body of work in affective computing shows that emotional state can be
estimated from facial expression and from autonomic signals such as heart-rate
variability, and organises emotion along the two dimensions of valence and
arousal~\cite{russell1980, posner2005}. Bringing such estimates to bear on music
selection is attractive, but two obstacles recur: continuously streaming camera
and physiological data to a server raises serious privacy concerns, and claims of
``mood-aware'' selection are rarely evaluated against a matched baseline.

We present \emph{Vibe Shuffle}, a client-only web application that closes this
loop entirely on-device. Facial expression is mapped to a continuous valence, an
optional Bluetooth heart-rate sensor contributes arousal, the two are fused onto
the valence--arousal circumplex, and the next track is chosen to match the fused
state from a fixed, curated pool. No audio, video, or physiological data leaves
the browser; only anonymous self-report ratings are exported.

Our contribution is twofold: a privacy-preserving, fully in-browser affect-to-music
pipeline, and a blinded, counterbalanced validation protocol that contrasts Vibe
Shuffle against a Random Shuffle baseline drawn from the same pool, isolating
perceived mood-fit from mere liking. The system, the curated pool, and the export
schema are openly available, making the design reproducible and falsifiable.

\section{Methodology}

\subsection{System overview}
Vibe Shuffle is a React/Vite single-page application with no backend. The signal
chain runs client-side: the camera feeds a facial valence estimator, an optional
heart-rate sensor feeds an arousal estimator, the channels are fused into a
$(\textit{valence}, \textit{arousal})$ state, that state is mapped to one of four
quadrants, and a ranking step selects the next track. Playback is handled by the
Spotify Web Playback SDK so that licensed audio plays in the browser without the
application ever receiving the audio stream.

\subsection{Valence from facial expression}
A MediaPipe Face Landmarker~\cite{mediapipe2019} produces per-frame blendshape
activations that are reduced to \emph{happy}, \emph{tense}, and \emph{sad}
scores. A slowly learned personal neutral baseline is subtracted to correct for
resting facial morphology, and the positive-versus-negative balance is mapped to a
continuous valence $v \in [0,1]$. Head and body motion (nose-tip drift and
frame-to-frame differencing) feed a motion channel that nudges the positive score
upward and contributes to arousal.

\subsection{Arousal from heart rate}
When a Bluetooth heart-rate sensor is connected, heart rate (HR) and inter-beat
(RR) intervals are parsed and a 120\,s personal baseline is built. Short-window
arousal combines the $z$-scored heart rate (higher HR raising arousal) with the
$z$-scored RMSSD, a standard short-term heart-rate-variability index whose
increase indicates lower arousal~\cite{shaffer2017}:
\begin{equation}
a_{\mathrm{phys}} = \operatorname{clip}\!\left(\tfrac{1}{2}
+ \alpha\, z(\mathrm{HR}) - \beta\, z(\mathrm{RMSSD}),\; 0,\; 1\right).
\end{equation}
SDNN is logged for signal-quality control but excluded from the short-window
estimate. The sensor is strictly optional; the system degrades gracefully without
it.

\subsection{Fusion and quadrants}
The fused state takes valence from the face and arousal primarily from the heart,
with head motion $m$ added on top:
\begin{equation}
(v, a) = \bigl(v_{\mathrm{face}},\; a_{\mathrm{phys}} + \gamma\, m\bigr).
\end{equation}
A usable electrocardiogram sets the arousal base in both directions; without one,
the face/motion channel carries arousal upward only; with neither face nor sensor,
both axes default to the neutral midpoint $0.5$. Splitting each axis at $0.5$
yields four quadrants---Energetic, Calm, Tense, and Melancholic---consistent with
the circumplex model of affect~\cite{russell1980, posner2005}.

\subsection{Music pool}
The pool is a fixed set of 100 tracks built once from the public Kaggle
\emph{Spotify Tracks Dataset} of real Spotify audio features~\cite{spotifydataset}.
Every track carries Spotify's two affect-bearing features on $[0,1]$:
\emph{valence}, the musical positivity of a track, and \emph{energy}, its
perceived intensity and activity. We take valence as the track's valence
coordinate and energy as its arousal coordinate, so each track $t$ inherits a
fixed $(v_t, a_t)$ with no per-session lookup. Tracks are categorised by the same
midpoint split used for the listener (\autoref{tab:quadrants}): \emph{Energetic}
when $v_t \ge 0.5$ and $a_t \ge 0.5$, \emph{Calm} when valence is high but energy
low, \emph{Tense} when valence is low but energy high, and \emph{Melancholic} when
both are low. To assemble a balanced pool, the most popular, de-duplicated tracks
with Spotify popularity $\ge 55$ were taken from each quadrant until 25 were
collected, yielding 100 tracks split evenly across the four quadrants. The same
pool is used for every participant, which keeps the Vibe-versus-Random comparison
controlled and removes any dependence on a participant's personal library.

\begin{table}[t]
\caption{Track categorisation. Each track's quadrant is fixed from its Spotify
valence $v_t$ and energy $a_t$ by splitting both axes at $0.5$; the pool holds 25
tracks per quadrant.}
\label{tab:quadrants}
\begin{tabularx}{\linewidth}{@{}lX@{}}
\toprule
Quadrant & Rule \\
\midrule
Energetic   & $v_t \ge 0.5,\ a_t \ge 0.5$ \\
Calm        & $v_t \ge 0.5,\ a_t < 0.5$ \\
Tense       & $v_t < 0.5,\ a_t \ge 0.5$ \\
Melancholic & $v_t < 0.5,\ a_t < 0.5$ \\
\bottomrule
\end{tabularx}
\end{table}

\subsection{Track selection}
Each track's quadrant is fixed in advance from its $(v_t, a_t)$, so selection
reduces to a draw over a candidate set. The fused listener state
$(v^\star, a^\star)$ is mapped to its quadrant $q^\star$. In a Vibe block the next
track is drawn from the mood-congruent quadrant only,
\begin{equation}
t^\star \in \{\, t : \operatorname{quad}(t) = q^\star \,\},
\end{equation}
with recently played tracks deprioritised so a block does not repeat songs. Within
that set the specific track is fixed by a per-session seed, making each run
reproducible yet different. A Random block applies the same seeded rule to the
entire pool. The two conditions thus share an identical pool and an identical
random mechanism; the \emph{only} manipulated variable is the quadrant constraint.
For transparency the Euclidean distance between the detected state and each track
is also computed and shown to the listener as a ``fit'' percentage, but it does
not determine which track plays.

\subsection{Validation protocol}
The application implements a blinded, counterbalanced within-subject comparison.
A session contains two blocks---\texttt{random} and \texttt{vibe}---of five tracks
each drawn from the shared pool of \autoref{tab:quadrants}; block order is
randomised per session and recorded for counterbalancing. Each track plays for 60\,s before the rating prompt opens, and
the condition is never revealed during the session. After each track, two 5-point
questions are asked in sequence: \emph{liking} (``How much do you like this
song?'') and \emph{mood-fit} (``How well did it fit your current mood?''). Mood-fit
is the primary outcome; liking is the control, separating ``did not fit my mood''
from ``I simply do not like this song.''

\section{Results}

This paper reports the implemented system and its pre-registered evaluation
design; a participant study using the protocol above is the immediate next step,
and no participant data is reported here. The working artifact demonstrates that
the full affect-to-music loop runs in real time, in-browser, on commodity
hardware with only a webcam and an optional consumer heart-rate strap.

Each completed session exports one row per track capturing the song features, the
detected $(v, a)$ state, signal-quality flags, and both ratings
(\autoref{tab:schema}). The pre-specified analysis contrasts the mood-fit rating
between \texttt{vibe} and \texttt{random} blocks, controlling for liking and
accounting for block order; liking is expected to be comparable across conditions
(both blocks draw from the same pool), while mood-fit is the discriminating
measure.

\begin{table}[t]
\caption{Selected exported fields (one row per played track).}
\label{tab:schema}
\begin{tabularx}{\linewidth}{@{}lX@{}}
\toprule
Field & Meaning \\
\midrule
\texttt{block\_mode} & condition: \texttt{vibe} or \texttt{random} \\
\texttt{song\_quadrant} & target affect quadrant of the track \\
\texttt{detected\_valence} & fused listener valence \\
\texttt{detected\_arousal} & fused listener arousal \\
\texttt{physiology\_quality} & heart-rate signal-quality flag \\
\texttt{rating\_fit\_1\_to\_5} & primary outcome: mood-fit \\
\texttt{rating\_like\_1\_to\_5} & control: liking \\
\bottomrule
\end{tabularx}
\end{table}

\section{Discussion}

By keeping the entire pipeline on-device, Vibe Shuffle sidesteps the privacy cost
that usually accompanies camera- and physiology-driven personalisation: raw frames
and heart-rate packets are processed in memory and discarded, and only the
anonymous rating CSV ever leaves the browser, at the participant's initiative.
This makes the system suitable for low-friction field deployment and for
classroom or seminar replication.

Several limitations bound the claims. The estimator is a transparent heuristic
rather than a validated clinical affect measure; facial valence is sensitive to
lighting and occlusion, and the arousal channel depends on a compatible sensor
that not every participant will own. The pool is small and fixed by design, so
results speak to relative ranking \emph{within} that pool rather than to
catalogue-scale recommendation. Finally, self-reported mood-fit is a perceptual
proxy; it captures whether selection \emph{feels} apt, which is the intended user
experience, but not an objective change in affective state.

\section{Conclusion}

Vibe Shuffle shows that a privacy-preserving, fully in-browser system can estimate
a listener's valence and arousal and select music to match it, and it pairs that
system with a blinded, counterbalanced protocol that can falsify the resulting
``mood-aware'' claim against a matched random baseline. The artifact and protocol
are openly available, providing a reproducible testbed; running the pre-registered
study to estimate the mood-fit advantage of state-aware selection is the natural
next step.
